\section*{Your turn}

Check out tag \code{ch12} of the companion repository and run \code{npm ci}
once. Exercises 1 to 5 need the database, the six subgraphs and the router.
Exercises 6 to 9 need node and nothing started at all: every one of them is a
composition.

Six services is more than a terminal each, so start them from the compose file.
The first run builds six images and is slow:

\begin{minted}[fontsize=\footnotesize,breaklines]{text}
docker compose up -d --build mosaic-db mosaic-catalog mosaic-pricing \
    mosaic-inventory mosaic-accounts mosaic-reviews mosaic-ordering mosaic-router
docker compose ps
\end{minted}

Eight containers, all healthy. If \code{mosaic-router} is the only one that is
not, port 3002 is taken, and the one other thing in this repository that wants
it is chapter~\ref{ch:how-a-federated-query-actually-runs}'s sample router:
\code{docker compose --profile wire down} clears it.
Chapter~\ref{ch:entity-resolution-done-right}'s router is behind the
\code{entities} profile and listens on 3003, so it is never the culprit.

\begin{enumerate}
  \item \textbf{Read the plan that grew.} Send this to port 3002 with the
  headers \code{X-WG-Include-Query-Plan: true} and
  \code{X-WG-Skip-Loader: true}:

  \begin{minted}[fontsize=\footnotesize,breaklines]{graphql}
{ browseProducts(first: 3) { nodes { title price { amount }
    reviews(first: 2) { nodes { author { displayName } } } } } }
  \end{minted}

  Find the \code{Parallel} node in \code{extensions.queryPlan} and count the
  fetches under it. There are two rather than the three this chapter printed,
  because this query asks Pricing and Reviews for something and Inventory for
  nothing. Then take \code{price} out and send it again. The node does not lose
  a child and carry on: it disappears, because one fetch is nothing to run in
  parallel and the planner stops describing it that way. Then take the reviews
  out too, and watch the accounts fetch go with them, because its keys came from
  the fetch you just deleted.

  \item \textbf{Find the fetch that could not be parallel.} In the first plan,
  read the \code{path} on the accounts fetch and the \code{dependsOnFetchIds}
  beside it. Write down, before looking, which fetch you expect it to depend on.
  The catalog fetch is the tempting answer, because that is where the query
  started and where every other dependency in this plan points.

  \item \textbf{Watch four services answer one question.} Send the same query
  without the two headers, so the loader runs. Then read the last timeline line
  out of each service:

  \begin{minted}[fontsize=\footnotesize,breaklines]{text}
docker compose logs --tail 5 mosaic-catalog mosaic-pricing mosaic-reviews mosaic-accounts
  \end{minted}

  Four services, four statements. Now ask the same of
  \code{mosaic-ordering} and notice that the answer is nothing.

  \item \textbf{Break a key format on purpose.} Open Pricing's copy of
  \code{ProductKey.cs}, under \code{Catalog/Model/}, and change the type name it
  checks from \code{nameof(Product)} to \code{"Widget"}. Rebuild Pricing and ask
  the router for a page of products with prices. Note which field is null, what
  the error says, and which service it names. Then work out which of the four
  copies of that file you would have had to check if the error had not named
  one. \code{git checkout} puts it back.

  \item \textbf{Submit a review from a customer who does not exist.} Take a
  product key out of Catalog, invent a \code{Customer} node identifier by
  base64-encoding \code{Customer:} followed by sixteen bytes of your choosing,
  and post \code{submitReview} to port 5105 with both. It succeeds. Then read
  the review back through the router and look at its author. Decide whether you
  would rather have had the write refused, and what it would have cost.

  \item \textbf{Make the composer say four different things about one
  directive.} Print these in turn:

  \begin{minted}[fontsize=\footnotesize,breaklines]{text}
node scripts/override-cases.mjs --print two-owners-no-override
node scripts/override-cases.mjs --print override-takes-the-field
node scripts/override-cases.mjs --print override-a-typo-with-a-real-loser
node scripts/override-cases.mjs --print override-from-yourself
  \end{minted}

  Two of them are the same error. Not one of the four contains the string
  \code{@override}; the closest any of them comes is the word \enquote{override}
  as a verb, in the case where you named your own subgraph, which is the mistake
  least likely to reach a build.

  \item \textbf{Suppress a warning and check nothing else moved.} Print
  \code{override-a-subgraph-that-is-not-there}, then read the case in
  \code{scripts/override-cases.mjs} to see how it proves that
  \code{--suppress-warnings} changes the output and not the config. Then change
  the case's expected routing so that it claims the field goes to
  \code{pricing}, and rerun to see what the failure tells you.

  \item \textbf{Hit the wall progressive \code{@override} is behind.} Print
  \code{progressive-override-is-rejected} and read the error. Then find
  \code{OverrideLegacySupportAttribute} in the HotChocolate source and satisfy
  yourself that the library would have emitted the label if the composer had
  asked for it. Two tools, both current, one specification version apart.

  \item \textbf{Override a key field, which composes.} Add
  \code{@override(from: "catalog")} to \code{Product.id} in a scratch copy of
  \code{schema/pricing.graphql} and compose the six by hand. It succeeds. Then
  read \code{engineConfig.datasourceConfigurations} in the output and find what
  catalog can no longer provide. This is the one thing in the lab the chapter
  itself does not cover: I composed it and read the routing table, and I never
  pointed a router at the result. If you do, you will know something this book
  does not.
\end{enumerate}

Exercise 4 is the one to do properly. Four services carry their own copy of a
file that decodes one key format, and nothing in any build would fail if one of
them drifted. The error you get is real and the question it raises is the one
worth carrying out of this chapter: the duplication that keeps five teams
independent is the same duplication that makes a format change a coordination
problem. Chapter~\ref{ch:ci-cd-for-a-federated-graph} is where a change like
that gets a workflow around it.
